\chapter{Stakeholders and User Stories}
\ChapterLabel{UserStories}
\index{stakeholder}
\index{user story}
\index{acceptance question}

\section{Chapter Responsibility}
\label{Section:UserStoriesResponsibility}
Identify the stakeholder roles, their desired capabilities and outcomes, and
the operating contexts and acceptance questions that guide validation.

\noindent\textbf{Expected result:} Four lecture-grounded stories with explicit
validation intent and visible decisions still needed for requirements.

\section{Source and Scope}
\label{Section:WeekTwoStoryScope}
Stories US-01--US-04 reproduce the Week 2 lecture wording verbatim
\cite[slides 19--20]{SPE800Week02}. The lecture identifies one role named
\emph{acquisition developer}. Additional stories will start at US-05;
these four IDs are reserved and will not be reassigned.
Operating contexts and acceptance questions below are proposed elaborations
for human review, not observed results or approved performance limits.

The lecture describes a sensor/ADC0804--FTDI245R--host acquisition path and
a mission application \cite[slides 12--15]{SPE800Week02}. This provides course
context; it does not approve project hardware or supersede the unresolved
choices in \ChapterRef{SystemContext} and
\hyperref[Table:WeekOneTBDs]{the Week 1 decision register}. Resolve those
choices before deriving binding requirements. The numerical teaching bounds
in the lecture are not adopted by these stories.

\section{Stakeholder Roles}
\label{Section:ADCStakeholders}
Roles describe beneficiaries and review perspectives, not assigned project
staff. Jack Marino coordinates stakeholder review under
\ChapterRef{DevelopmentPlan}; individual representatives are TBD.
\begin{tabularx}{\linewidth}{L{3.4cm}Y L{1.7cm}}
\caption{Stakeholder Roles\label{Table:ADCStakeholders}}\\
\toprule
\textbf{Role} & \textbf{Stakeholder concern} & \textbf{Story}\\
\midrule
\endfirsthead
\toprule
\textbf{Role} & \textbf{Stakeholder concern} & \textbf{Story}\\
\midrule
\endhead
\bottomrule
\endlastfoot
Mission operator & Decide whether sensor readings are usable through the mission application. & \UserStoryRef{US-01}\\
Acquisition developer & Sustain acquisition while processing and delivery proceed; understand degraded streaming. & \UserStoryRef{US-02}\\
Field engineer & Deploy and operate the subsystem on available field computers. & \UserStoryRef{US-03}\\
Mission integrator & Interpret a documented interface and replace compatible sensors or consumers. & \UserStoryRef{US-04}\\
\end{tabularx}

\section{User Story Register}
\label{Section:ADCUserStories}
All four stories are candidates for the course scope. Their proposed priority
is \gls{must}, subject to owner and course review. Requirements and use cases
will be allocated in their authoritative chapters; generic service-request
requirements do not implement these ADC stories. The
\hyperref[Table:TraceabilityMatrix]{Traceability Register} exposes the pending
mappings and links each story to a planned validation activity. A positive
answer to an acceptance question requires stakeholder observations in the
agreed context; document compilation is not stakeholder validation.

\begin{tabularx}{\linewidth}{L{1.6cm}Y L{2.3cm}}
\caption{User Story Register\label{Table:UserStories}}\\
\toprule
\textbf{ID} & \textbf{Role} & \textbf{Validation plan}\\
\midrule
\endfirsthead
\toprule
\textbf{ID} & \textbf{Role} & \textbf{Validation plan}\\
\midrule
\endhead
\bottomrule
\endlastfoot
\UserStoryRef{US-01} & Mission Operator & \ValidationRef{VAL-01}\\
\UserStoryRef{US-02} & Acquisition Developer & \ValidationRef{VAL-02}\\
\UserStoryRef{US-03} & Field Engineer & \ValidationRef{VAL-03}\\
\UserStoryRef{US-04} & Mission Integrator & \ValidationRef{VAL-04}\\
\end{tabularx}

\subsection{US-01: Mission Operator}
\noindent\textbf{ID:} \UserStoryID{US-01}\quad\textbf{Priority:} \gls{must} (proposed).

\noindent\textbf{Story.} As a mission operator, I need to see current sensor values and their availability, validity, and age so that I can decide whether to use a reading.

\paragraph{Operating context.} During a representative mission workflow, the mission operator views readings through the mission application. The scenario includes a current valid reading, interrupted input, stale output, and an unavailable sensor. The display and field conditions remain to be agreed with the stakeholder.

\paragraph{Acceptance questions.}
\begin{enumerate}
\item Can the mission operator identify which reading is current and valid and decide whether to use it?
\item When input stops or the sensor becomes unavailable, can the mission operator distinguish stale or unavailable data from usable data?
\item Are availability, validity, and age understandable without coaching in the agreed operating context?
\end{enumerate}
\paragraph{Open decisions.} Agree representative field conditions, sample-age meaning, stale-state criteria, and expected use/no-use decisions. Timing allocations remain TBD.

\noindent\textbf{Validation:} \ValidationRef{VAL-01} (planned; \StatusPending). Requirements and use-case allocation: TBD before requirements baseline.

\subsection{US-02: Acquisition Developer}
\noindent\textbf{ID:} \UserStoryID{US-02}\quad\textbf{Priority:} \gls{must} (proposed).

\noindent\textbf{Story.} As an acquisition developer, I need sensor acquisition to continue while prior samples are processed and delivered so that mission software receives a continuous stream of readings.

\paragraph{Operating context.} In the lecture scenario, the sensor and ADC0804 supply an 8-bit code through an FTDI245R USB interface. USB reads, processing, and output writes proceed concurrently on the host. Review normal load, a slow consumer, overload, and disconnect/recovery; hardware selection and load bounds remain project decisions.

\paragraph{Acceptance questions.}
\begin{enumerate}
\item Can the acquisition developer show that new reads continue while earlier samples are processed and delivered at the agreed operating load?
\item Can the acquisition developer determine whether samples were delayed, lost, or invalid when processing or output falls behind?
\item Can mission software interpret the stream and its status through disconnect and recovery without mistaking stale data for new data?
\end{enumerate}
\paragraph{Open decisions.} Define sample rate, supported processing load, queue capacity, loss policy, timing bounds, clock, and recovery behavior before validation. Continuous streaming does not establish an unconditional zero-loss guarantee.

\noindent\textbf{Validation:} \ValidationRef{VAL-02} (planned; \StatusPending). Requirements and use-case allocation: TBD before requirements baseline.

\subsection{US-03: Field Engineer}
\noindent\textbf{ID:} \UserStoryID{US-03}\quad\textbf{Priority:} \gls{must} (proposed).

\noindent\textbf{Story.} As a field engineer, I need the sensor subsystem to run on macOS, Windows, and Linux so that I can use the computers available in the field.

\paragraph{Operating context.} A field engineer installs, configures, and runs the same sensor-subsystem capability on representative macOS, Windows, and Linux field computers. OS versions, processor architectures, FTDI driver dependencies, permissions, and available resources must be selected before evaluating portability.

\paragraph{Acceptance questions.}
\begin{enumerate}
\item Can the field engineer install and start the subsystem on each agreed macOS, Windows, and Linux configuration using the documented procedure?
\item Can the field engineer acquire readings and inspect status on each platform without platform-specific source changes?
\item Can the field engineer recognize and resolve a missing driver, permission, or configuration dependency using the deployment instructions?
\end{enumerate}
\paragraph{Open decisions.} Agree the supported platform matrix, installation prerequisites, deployment procedure, and representative field constraints. No platform support has yet been demonstrated.

\noindent\textbf{Validation:} \ValidationRef{VAL-03} (planned; \StatusPending). Requirements and use-case allocation: TBD before requirements baseline.

\subsection{US-04: Mission Integrator}
\noindent\textbf{ID:} \UserStoryID{US-04}\quad\textbf{Priority:} \gls{must} (proposed).

\noindent\textbf{Story.} As a mission integrator, I need a documented sensor data and status interface so that I can replace a sensor or mission consumer without redesigning their connection.

\paragraph{Operating context.} A mission integrator uses a documented contract between the sensor subsystem and mission software. Local sample records and the mission-facing interface are separate boundaries. The lecture specifies OMS for the course mission interface; program applicability, standard version, CAL profile, UCI mapping, and service ownership remain open.

\paragraph{Acceptance questions.}
\begin{enumerate}
\item Can the mission integrator interpret sensor data and availability, validity, age, and fault status from the documented contract without undocumented assumptions?
\item Can a contract-compatible replacement sensor or mission consumer be integrated without redesigning the agreed connection?
\item Are local record semantics and the proposed OMS mission exchange distinguished, with unresolved profile and mapping decisions visible?
\end{enumerate}
\paragraph{Open decisions.} Agree the interface contract, compatibility boundary, data/status mapping, and replacement scenario. Confirm OMS applicability, version, CAL profile, UCI mapping, and gateway ownership before conformance claims.

\noindent\textbf{Validation:} \ValidationRef{VAL-04} (planned; \StatusPending). Requirements and use-case allocation: TBD before requirements baseline.

\section{Review Rules}
Check each story with its stakeholder for role, capability, outcome, operating
context, and acceptance questions. Confirm priority and settle the open
decisions before baselining requirements. Add reciprocal requirement and
use-case links only when those artifacts exist; preserve the story-to-validation
links and record observed stakeholder results in \ChapterRef{Validation}.

\section{Retired Template Examples}
\label{Section:RetiredTemplateStories}
The following service-request examples are retained only to preserve template
IDs and links in chapters awaiting tailoring. They are not ADC stakeholder
needs, approved requirements, or validation evidence. IDs US-001--US-003 are
distinct from the lecture IDs US-01--US-04 and will not be reused.
\begin{tabularx}{\linewidth}{L{1.9cm}YL{1.8cm}L{3.4cm}}
\caption{Retired Template Story Register\label{Table:TemplateUserStories}}\\
\toprule
\textbf{ID} & \textbf{Story and Acceptance Criteria} & \textbf{Priority} & \textbf{Requirement(s)}\\
\midrule
\endfirsthead
\multicolumn{4}{l}{\small\itshape Retired Template Story Register (continued)}\\
\toprule
\textbf{ID} & \textbf{Story and Acceptance Criteria} & \textbf{Priority} & \textbf{Requirement(s)}\\
\midrule
\endhead
\midrule
\multicolumn{4}{r}{\small\itshape Continued on next page}\\
\endfoot
\bottomrule
\endlastfoot
\UserStoryID{US-001} & As a customer, I want to submit a service request so that the responsible team can act on it. Acceptance: valid data produces a unique request identifier on every supported client. & \gls{must} & \RequirementRef{REQ-F-001}, \RequirementRef{REQ-C-001}\\
\UserStoryID{US-002} & As a customer, I want to view my request status so that I understand its progress. Acceptance: the current status and most recent update appear within the specified response time, and the transition history is retained. & \gls{should} & \RequirementRef{REQ-F-002}, \RequirementRef{REQ-Q-001}, \RequirementRef{REQ-D-001}, \RequirementRef{REQ-C-001}\\
\UserStoryID{US-003} & As an operations analyst, I want to export request history so that I can perform approved analysis. Acceptance: an authorized actor prepares complete history within five minutes in the approved versioned format without changing source records. & \gls{would} & \RequirementRef{REQ-B-001}, \RequirementRef{REQ-F-003}, \RequirementRef{REQ-I-001}, \RequirementRef{REQ-S-001}, \RequirementRef{REQ-D-001}, \RequirementRef{REQ-C-001}\\
\end{tabularx}
